\documentclass[12pt,a4paper]{article}

% ── Packages ──────────────────────────────────────────────────────────────────
\usepackage[T1]{fontenc}
\usepackage[utf8]{inputenc}
\usepackage{lmodern}
\usepackage[margin=2.5cm]{geometry}
\usepackage{amsmath,amssymb,amsthm}
\usepackage{mathtools}
\usepackage{booktabs}
\usepackage{longtable}
\usepackage{array}
\usepackage{xcolor}
\usepackage{hyperref}
\usepackage{enumitem}
\usepackage{parskip}
\usepackage{microtype}
\usepackage{graphicx}
\usepackage{float}
\usepackage{caption}
\usepackage{subcaption}
\usepackage{algorithm}
\usepackage{algpseudocode}
\usepackage{listings}
\usepackage{tcolorbox}
\usepackage{mdframed}
\usepackage{titlesec}
\usepackage{fancyhdr}
\usepackage{setspace}

% ── Colours ───────────────────────────────────────────────────────────────────
\definecolor{clipblue}{RGB}{70,114,176}
\definecolor{tdaorange}{RGB}{221,132,82}
\definecolor{fttagreen}{RGB}{85,168,104}
\definecolor{lightgray}{RGB}{245,245,245}
\definecolor{darkgray}{RGB}{80,80,80}
\definecolor{linkblue}{RGB}{0,70,180}

% ── Hyperref ──────────────────────────────────────────────────────────────────
\hypersetup{
  colorlinks=true,
  linkcolor=linkblue,
  citecolor=linkblue,
  urlcolor=linkblue,
  pdftitle={CLIP, Domain Shift, and Test-Time Adaptation: TDA vs FreeTTA},
  pdfauthor={Research Report},
}

% ── Header / Footer ───────────────────────────────────────────────────────────
\pagestyle{fancy}
\fancyhf{}
\fancyhead[L]{\small CLIP Test-Time Adaptation: TDA vs FreeTTA}
\fancyhead[R]{\small\thepage}
\renewcommand{\headrulewidth}{0.4pt}

% ── Section formatting ────────────────────────────────────────────────────────
\titleformat{\section}{\large\bfseries\color{darkgray}}{\thesection.}{0.5em}{}[\titlerule]
\titleformat{\subsection}{\normalsize\bfseries\color{darkgray}}{\thesubsection.}{0.5em}{}
\titleformat{\subsubsection}{\normalsize\bfseries}{\thesubsubsection.}{0.5em}{}

% ── Math operators ────────────────────────────────────────────────────────────
\DeclareMathOperator*{\argmax}{arg\,max}
\DeclareMathOperator*{\argmin}{arg\,min}
\newcommand{\R}{\mathbb{R}}
\newcommand{\bx}{\mathbf{x}}
\newcommand{\bw}{\mathbf{w}}
\newcommand{\bW}{\mathbf{W}}
\newcommand{\bz}{\mathbf{z}}
\newcommand{\bmu}{\boldsymbol{\mu}}
\newcommand{\1}{\mathbf{1}}
\newcommand{\norm}[1]{\left\|#1\right\|}

% ── tcolorbox styles ──────────────────────────────────────────────────────────
\tcbset{
  intuitionbox/.style={
    colback=lightgray, colframe=darkgray!50,
    boxrule=0.5pt, arc=3pt,
    left=6pt, right=6pt, top=4pt, bottom=4pt,
    fonttitle=\bfseries\small, title={Intuition},
  },
  findingbox/.style={
    colback=fttagreen!8, colframe=fttagreen!60,
    boxrule=0.8pt, arc=3pt,
    left=6pt, right=6pt, top=4pt, bottom=4pt,
    fonttitle=\bfseries\small,
  },
}

% ── Listings (code blocks) ────────────────────────────────────────────────────
\lstset{
  backgroundcolor=\color{lightgray},
  basicstyle=\ttfamily\footnotesize,
  breaklines=true,
  frame=single,
  framerule=0.4pt,
  rulecolor=\color{darkgray!30},
}

% ── Document ──────────────────────────────────────────────────────────────────
\begin{document}

% ── Title Page ────────────────────────────────────────────────────────────────
\begin{titlepage}
  \centering
  \vspace*{2cm}
  {\LARGE\bfseries Comprehensive Research Report\par}
  \vspace{0.5cm}
  {\Large CLIP, Domain Shift, and Test-Time Adaptation:\par}
  {\Large TDA vs FreeTTA --- A Deep Comparative Study\par}
  \vspace{1cm}
  \rule{0.7\textwidth}{0.4pt}
  \vspace{0.5cm}

  {\normalsize
  \textit{Datasets:} Caltech101 $\cdot$ DTD $\cdot$ EuroSAT $\cdot$ ImageNet-V2 $\cdot$ Oxford Pets\\[4pt]
  \textit{Pipeline:} \texttt{experiments/run\_comparative\_analysis.py}\\[4pt]
  \textit{Stream order:} Natural $\cdot$ Rolling window $= 50$ $\cdot$ Seed $= 42$
  }

  \vspace{1cm}
  \rule{0.7\textwidth}{0.4pt}
  \vspace{2cm}

  {\small All numerical results from \texttt{outputs/comparative\_analysis/}.
  Plots available in the same directory.}
\end{titlepage}

% ── Table of Contents ─────────────────────────────────────────────────────────
\tableofcontents
\clearpage

% ══════════════════════════════════════════════════════════════════════════════
\section{CLIP: What, How, and Why}
% ══════════════════════════════════════════════════════════════════════════════

\subsection{What is CLIP?}

CLIP (Contrastive Language--Image Pretraining) is a vision--language model developed by OpenAI \cite{radford2021clip}.
Its defining feature is a \textbf{shared embedding space} for both images and text, learned by training on 400~million image--caption pairs scraped from the internet.
Once trained, CLIP can match any image to any text description---including class names it has never explicitly seen.

This gives CLIP its most remarkable property: \textbf{zero-shot classification}.
Given a set of class names defined at inference time, CLIP can classify images with no labelled training examples, simply by comparing image features to text-encoded class descriptions.

\subsection{Architecture}

CLIP consists of two parallel encoders trained jointly:

\begin{itemize}[leftmargin=1.5em]
  \item \textbf{Image encoder} $f_I$: Vision Transformer ViT-B/16.
        Takes a $224\times224$ image and outputs a 512-dimensional unit vector.
  \item \textbf{Text encoder} $f_T$: Transformer.
        Takes a text string (e.g.\ \textit{``a photo of a cat''}) and outputs a 512-dimensional unit vector.
\end{itemize}

Both encoders project to the same 512-D space, where semantic similarity translates to geometric proximity (high cosine similarity).

\subsection{Training: Contrastive Loss}

For a batch of $N$ image--text pairs $\{(I_i, t_i)\}_{i=1}^N$, CLIP maximises the cosine similarity of matched pairs and minimises it for all $N^2 - N$ mismatched pairs.
The loss is symmetric cross-entropy applied to an $N\times N$ similarity matrix:

\begin{equation}
  s(I_i, t_j) = f_I(I_i)^\top f_T(t_j)
  \qquad \text{(both unit-normed, so this equals cosine similarity)}
  \label{eq:cosine}
\end{equation}

\begin{equation}
  \mathcal{L}_{\text{img}} = -\frac{1}{N}\sum_{i=1}^N \log
    \frac{\exp(\tau \cdot s(I_i,t_i))}{\sum_{j=1}^N \exp(\tau \cdot s(I_i,t_j))}
  \label{eq:loss_img}
\end{equation}

\begin{equation}
  \mathcal{L}_{\text{txt}} = -\frac{1}{N}\sum_{j=1}^N \log
    \frac{\exp(\tau \cdot s(I_j,t_j))}{\sum_{i=1}^N \exp(\tau \cdot s(I_i,t_j))}
  \label{eq:loss_txt}
\end{equation}

\begin{equation}
  \mathcal{L} = \frac{\mathcal{L}_{\text{img}} + \mathcal{L}_{\text{txt}}}{2}
  \label{eq:loss_total}
\end{equation}

where $\tau$ is a learned temperature parameter.

\begin{tcolorbox}[intuitionbox]
Pull the image and its paired caption together; push every other image--caption combination apart.
After 400M examples, the shared space encodes rich cross-modal semantics.
\end{tcolorbox}

\subsection{Zero-Shot Inference}

To classify a query image $I$ into $C$ classes:

\begin{enumerate}[leftmargin=1.5em]
  \item Build text templates: for each class $c$, form prompt $t_c = $ \textit{``a photo of a \{class\_name\}''}.
  \item Encode all $C$ prompts:
    $\bW = [f_T(t_1), \ldots, f_T(t_C)] \in \R^{C \times D}$, each row unit-normed.
  \item Encode query image: $\bx = f_I(I) \in \R^D$, unit-normed.
  \item Compute logits:
    \begin{equation}
      \bz = 100 \cdot \bW \bx \in \R^C
      \quad \text{(scale factor 100 $\approx 1/\tau$)}
      \label{eq:clip_logits}
    \end{equation}
  \item Probabilities:
    \begin{equation}
      p_c = \mathrm{softmax}(\bz)_c = \frac{\exp(z_c)}{\sum_{j=1}^C \exp(z_j)}
      \label{eq:softmax}
    \end{equation}
  \item Predict: $\hat{y} = \argmax_c p_c$
\end{enumerate}

\subsection{Strengths and Weaknesses}

\begin{minipage}[t]{0.48\textwidth}
\textbf{Strengths}
\begin{itemize}[leftmargin=1em,noitemsep]
  \item Zero-shot: no labelled data needed
  \item Single frozen model covers hundreds of tasks
  \item Rich semantics from 400M training pairs
\end{itemize}
\end{minipage}\hfill
\begin{minipage}[t]{0.48\textwidth}
\textbf{Weaknesses}
\begin{itemize}[leftmargin=1em,noitemsep]
  \item Satellite/medical/microscopy images underrepresented
  \item Template ``a photo of'' mismatches non-photographic domains
  \item Near-uniform confidence: $p_c \approx 1/C$ for all $c$
\end{itemize}
\end{minipage}

% ══════════════════════════════════════════════════════════════════════════════
\section{The Problem: Domain Shift}
% ══════════════════════════════════════════════════════════════════════════════

\subsection{Definition}

Domain shift means the statistical distribution of test images $P_{\text{test}}(\bx)$ differs from the distribution the model was trained on $P_{\text{train}}(\bx)$, even if the label set is identical.
Formally:
\begin{equation}
  P_{\text{train}}(\bx \mid y) \;\neq\; P_{\text{test}}(\bx \mid y)
  \label{eq:domain_shift}
\end{equation}

\subsection{Why Domain Shift Hurts CLIP Specifically}

CLIP's text anchors $\bW$ are \emph{fixed} at inference time---they encode the ``average internet concept'' of each class.
When test images deviate from this concept, cosine similarities shift and the wrong class becomes top-ranked.

\begin{description}
  \item[\textbf{EuroSAT:}] The anchor for ``Annual Crop'' encodes a landscape photograph perspective.
  Actual EuroSAT images are $64\times64$ multispectral overhead tiles.
  The feature vector $\bx$ lands equidistant from multiple class anchors, yielding highly uncertain predictions.

  \item[\textbf{DTD:}] Texture names like ``braided'', ``woven'', and ``knitted'' are semantically similar in language.
  CLIP's text encoder clusters these anchors close together, making purely text-based discrimination hard.
\end{description}

\subsection{Quantifying Domain Shift via CLIP Baseline Accuracy}

\begin{table}[H]
\centering
\caption{CLIP zero-shot accuracy and dataset characteristics.}
\label{tab:clip_baseline}
\begin{tabular}{@{}lrrrl@{}}
\toprule
Dataset & Classes & Samples & CLIP Acc.\ (\%) & Domain Type \\
\midrule
Caltech101  & 100   & 2,465  & 93.55 & Standard photographs \\
Oxford Pets & 37    & 3,669  & 88.39 & Fine-grained pet breeds \\
ImageNet-V2 & 1,000 & 10,000 & 62.35 & General objects (shifted test set) \\
DTD         & 47    & 1,880  & 43.94 & Abstract textures \\
EuroSAT     & 10    & 8,100  & 48.43 & Satellite remote sensing \\
\bottomrule
\end{tabular}
\end{table}

The 44.6-point gap between Caltech (93.55\%) and EuroSAT (48.43\%) directly measures domain shift severity.
On EuroSAT, CLIP barely exceeds random chance for a 10-class problem ($1/10 = 10\%$).

\subsection{CLIP Confidence Is Uninformative}

\begin{equation}
  \text{Expected uniform confidence for $C$ classes} = \frac{1}{C}
\end{equation}

Observed: Caltech mean confidence $= 0.0113$ (uniform $\approx 0.010$ for 100 classes); EuroSAT mean confidence $= 0.103$ (uniform $= 0.100$ for 10 classes).
CLIP cannot distinguish easy from hard samples.

% ══════════════════════════════════════════════════════════════════════════════
\section{Test-Time Adaptation: The Solution}
% ══════════════════════════════════════════════════════════════════════════════

\subsection{What is TTA?}

Test-Time Adaptation (TTA) adapts a frozen model at inference time using only the unlabelled test stream.
Constraints:
\begin{itemize}[noitemsep]
  \item Access only to the current and past test samples (online)
  \item CLIP weights remain completely frozen
  \item Decision on sample $t$ is made before seeing sample $t+1$
  \item No ground-truth labels: adaptation signal comes from the model's own predictions
\end{itemize}

\subsection{Why TTA Works}

Even without labels, the test stream carries three types of signal:
\begin{enumerate}[leftmargin=1.5em,noitemsep]
  \item \textbf{Self-consistency}: samples that look similar to each other and receive the same high-confidence prediction probably share a true class label.
  \item \textbf{Cluster structure}: in a well-trained feature space, same-class samples cluster; TTA methods exploit this geometry to shift decision boundaries.
  \item \textbf{Entropy minimisation}: confident predictions (low entropy) are statistically more likely correct; updating class representations from them improves future predictions.
\end{enumerate}

\subsection{Why Not Fine-Tune?}

CLIP's visual encoder has 307M parameters.
Fine-tuning on a few hundred test samples causes catastrophic forgetting of general knowledge.
TTA instead adapts in the downstream logit space---a much lower-dimensional, well-conditioned problem.

% ══════════════════════════════════════════════════════════════════════════════
\section{Motivation: Why TDA and FreeTTA?}
% ══════════════════════════════════════════════════════════════════════════════

\subsection{Two Competing Philosophies}

After surveying TTA methods for CLIP, two fundamentally different philosophies emerged:

\begin{tcolorbox}[findingbox, title={Philosophy 1 --- ``Remember What You've Seen'' (TDA)}]
  Keep a running memory bank of reliable past predictions.
  When a new sample arrives, use its similarity to cached samples to refine the prediction.
  \textbf{Local}: each prediction is influenced only by previously seen, similar-looking samples.
\end{tcolorbox}

\vspace{0.4em}

\begin{tcolorbox}[findingbox, title={Philosophy 2 --- ``Update Your World Model'' (FreeTTA)}]
  Maintain running estimates of each class's feature distribution (class means).
  Predictions fuse CLIP's text anchor with the learned visual mean.
  Update means after each prediction.
  \textbf{Global}: each prediction is influenced by everything seen for that class so far.
\end{tcolorbox}

\subsection{Why These Two Specifically?}

\begin{enumerate}[leftmargin=1.5em]
  \item \textbf{Minimal assumptions}: neither requires labelled data, neither modifies CLIP weights.
  \item \textbf{Complementary mechanisms}: cache is instance-level; EM is class-level.
        Comparing them reveals which granularity of signal matters.
  \item \textbf{Different failure modes}: cache fails under memory pressure; EM fails with few samples per class.
        Each failure exposes the other's strength.
  \item \textbf{Practical efficiency}: both run in $O(1)$ or $O(C)$ per step with no backpropagation.
  \item \textbf{Theoretical grounding}: TDA relates to $k$-NN retrieval; FreeTTA to Gaussian mixture EM.
\end{enumerate}

% ══════════════════════════════════════════════════════════════════════════════
\section{TDA: Cache-Based Adaptation}
% ══════════════════════════════════════════════════════════════════════════════

\subsection{Intuition}

\begin{tcolorbox}[intuitionbox]
Imagine a new doctor who, in addition to a textbook (CLIP's text anchors), keeps notes on recent patients.
For each new patient, they consult the 5 most similar past cases.
If those notes confidently indicate ``Annual Crop'', the doctor upweighs that diagnosis---even if the textbook description of Annual Crop doesn't perfectly match satellite imagery.
\end{tcolorbox}

TDA maintains two cache types per class:
\begin{itemize}
  \item \textbf{Positive cache}: recent samples with low prediction entropy (high confidence). Good reference examples.
  \item \textbf{Negative cache}: samples with medium entropy. Anti-examples---``this looks somewhat like class $c$ but probably isn't.''
\end{itemize}

\subsection{Algorithm}

\begin{algorithm}[H]
\caption{TDA Online Adaptation}\label{alg:tda}
\begin{algorithmic}[1]
\State \textbf{Initialise:} $\mathcal{P}_c \leftarrow \emptyset$,\; $\mathcal{N}_c \leftarrow \emptyset$ for all $c \in [C]$
\For{each test sample $\bx_t$ (unit-normed $\in\R^D$)}
  \State \textbf{// Step 1: CLIP logits}
  \State $z^{\text{clip}}_c = 100 \cdot \bx_t^\top \bw_c$ \quad for all $c$
  \State \textbf{// Step 2: Cache adjustment logits}
  \State $z^{\text{pos}}_c = \dfrac{\lambda_+}{|\mathcal{P}_c|}\displaystyle\sum_{\mathbf{v}\in\mathcal{P}_c} \bx_t^\top \mathbf{v}$ \quad (0 if $\mathcal{P}_c = \emptyset$)
  \State $z^{\text{neg}}_c = \dfrac{\lambda_-}{|\mathcal{N}_c|}\displaystyle\sum_{\mathbf{v}\in\mathcal{N}_c} \bx_t^\top \mathbf{v}$ \quad (0 if $\mathcal{N}_c = \emptyset$)
  \State \textbf{// Step 3: Combine and predict}
  \State $z^{\text{final}}_c = z^{\text{clip}}_c + z^{\text{pos}}_c - z^{\text{neg}}_c$
  \State $p(c \mid \bx_t) = \mathrm{softmax}(\mathbf{z}^{\text{final}})_c$
  \State $\hat{y} \leftarrow \argmax_c p(c \mid \bx_t)$
  \State \textbf{// Step 4: Update cache}
  \State $H \leftarrow -\sum_c p(c|\bx_t)\log p(c|\bx_t)$
  \If{$H < \theta_{\text{low}}$} add $\bx_t$ to $\mathcal{P}_{\hat{y}}$; evict oldest if over capacity
  \ElsIf{$H < \theta_{\text{mid}}$} add $\bx_t$ to $\mathcal{N}_{\hat{y}}$; evict oldest if over capacity
  \EndIf
\EndFor
\end{algorithmic}
\end{algorithm}

\subsection{Mathematics}

The positive cache logit is a weighted cosine similarity kernel:
\begin{equation}
  z^{\text{pos}}_c(\bx_t) = \lambda_+ \cdot \mathbb{E}_{\mathbf{v} \sim \mathcal{P}_c}\!\left[\bx_t^\top \mathbf{v}\right]
  = \lambda_+ \cdot \bx_t^\top \hat{\bmu}_c
  \label{eq:tda_kernel}
\end{equation}
where $\hat{\bmu}_c = \frac{1}{|\mathcal{P}_c|}\sum_{\mathbf{v}\in\mathcal{P}_c}\mathbf{v}$ is the empirical centroid of cached features.

In the limit of infinite cache capacity, $\hat{\bmu}_c \to \bmu_c$ (true class centroid), making TDA equivalent to a nearest-centroid classifier.
With finite capacity $K$ (default $K=3$), it is a \textbf{sliding-window mean of the $K$ most recent confident examples}.

The combined logit adjustment is a kernel regression:
\begin{equation}
  \mathbf{z}^{\text{final}} = \mathbf{z}^{\text{clip}} + \sum_{c} \bigl[\lambda_+ K^+(\bx_t, c) - \lambda_- K^-(\bx_t, c)\bigr] \mathbf{e}_c
  \label{eq:tda_combined}
\end{equation}
where $K^+$ and $K^-$ are dot-product similarity kernels to positive and negative caches respectively.

\subsection{Hyperparameters}

\begin{table}[H]
\centering
\caption{TDA hyperparameters (uniform across all datasets in our experiments).}
\label{tab:tda_hyperparams}
\begin{tabular}{@{}llll@{}}
\toprule
Parameter & Symbol & Value & Meaning \\
\midrule
Positive shot capacity & $K^+$ & 3 & Max positive samples per class \\
Negative shot capacity & $K^-$ & 3 & Max negative samples per class \\
Low entropy threshold  & $\theta_{\text{low}}$ & 0.20 & Entry threshold for positive cache \\
Mid entropy threshold  & $\theta_{\text{mid}}$ & 0.40 & Entry threshold for negative cache \\
Positive cache weight  & $\lambda_+$ & 1.0 & Weight on positive adjustment \\
Negative cache weight  & $\lambda_-$ & 0.1 & Weight on negative adjustment \\
\bottomrule
\end{tabular}
\end{table}

\subsection{Key Structural Properties}

\begin{itemize}
  \item \textbf{Local}: each prediction depends only on the most recently seen similar samples (FIFO window).
  \item \textbf{Cold-start}: no adjustment until cache fills; first steps are essentially raw CLIP.
  \item \textbf{Cache pressure}: when $N_{\text{total}} \gg C \times K^+$, slots are overwritten faster than useful signal accumulates.
  \item \textbf{Discriminative}: negative cache actively suppresses confusion between visually similar classes.
  \item \textbf{No global state}: two samples from the same class seen 1,000 steps apart do not interact.
\end{itemize}

% ══════════════════════════════════════════════════════════════════════════════
\section{FreeTTA: Online EM Adaptation}
% ══════════════════════════════════════════════════════════════════════════════

\subsection{Intuition}

\begin{tcolorbox}[intuitionbox]
The same new doctor, but instead of notes on individual patients, they build a running average of what each disease ``looks like'' in test images.
Every time ``Annual Crop'' is diagnosed, the visual prototype for Annual Crop shifts slightly toward the current image.
Over 800 samples, the prototype converges to the true distribution of satellite-image crops---even if the textbook description was wrong.
\end{tcolorbox}

FreeTTA models each class as a Gaussian distribution in feature space and adapts class means online via the Expectation-Maximisation (EM) framework.

\subsection{Algorithm}

\begin{algorithm}[H]
\caption{FreeTTA Online EM Adaptation}\label{alg:freetta}
\begin{algorithmic}[1]
\State \textbf{Initialise:} $\bmu_c \leftarrow \bw_c$ for all $c$ \quad (CLIP text anchor)
\For{each test sample $\bx_t$ (unit-normed)}
  \State \textbf{// E-Step: compute adapted probabilities}
  \State $z^{\text{clip}}_c = 100 \cdot \bx_t^\top \bw_c$
  \State $\alpha_t = 1 - H_{\text{clip}}(\bx_t) / \log C$ \quad (EM weight from CLIP entropy)
  \State $z^{\text{adapt}}_c = \alpha_t \cdot \bx_t^\top \bmu_c / \tau_{\text{adapt}}$
  \State $p_{\text{adapt}}(c \mid \bx_t) = \mathrm{softmax}\!\left(\mathbf{z}^{\text{clip}} + \mathbf{z}^{\text{adapt}}\right)_c$
  \State $\hat{y} \leftarrow \argmax_c\, p_{\text{adapt}}(c \mid \bx_t)$
  \State \textbf{// M-Step: update predicted class mean}
  \State $w_{\text{upd}} = \sigma\!\left(\gamma \cdot (p_{\text{adapt}}(\hat{y} \mid \bx_t) - 0.5)\right)$
  \State $\bmu_{\hat{y}} \leftarrow (1 - \eta \cdot w_{\text{upd}})\,\bmu_{\hat{y}} + \eta \cdot w_{\text{upd}}\,\bx_t$
  \State $\bmu_{\hat{y}} \leftarrow \bmu_{\hat{y}} / \norm{\bmu_{\hat{y}}}$ \quad (re-normalise)
\EndFor
\end{algorithmic}
\end{algorithm}

\subsection{Mathematics}

\subsubsection{The EM View}

FreeTTA maximises the expected log-likelihood of test data under a Gaussian mixture model:
\begin{equation}
  p(\bx) = \sum_{c=1}^C \pi_c \cdot \mathcal{N}(\bx;\, \bmu_c,\, \sigma^2 \mathbf{I})
  \label{eq:gmm}
\end{equation}
where $\pi_c = p_{\text{clip}}(c)$ (CLIP prior) and covariances are isotropic.

The \textbf{E-step} assigns soft responsibilities:
\begin{equation}
  r_c(\bx_t) = p_{\text{adapt}}(c \mid \bx_t)
\end{equation}

The \textbf{M-step} performs stochastic gradient ascent on the mixture log-likelihood:
\begin{equation}
  \frac{\partial}{\partial \bmu_c} \log p(\bx_t) \approx \frac{r_c(\bx_t)\,(\bx_t - \bmu_c)}{\sigma^2}
  \label{eq:mstep_grad}
\end{equation}
\begin{equation}
  \bmu_c \leftarrow \bmu_c + \eta \cdot r_c(\bx_t) \cdot (\bx_t - \bmu_c)
  \label{eq:online_kmeans}
\end{equation}
This is exactly the \textbf{online $k$-means update}, weighted by soft responsibility.

\subsubsection{The Drift Equation}

Over $T$ updates to class $c$, the mean follows an exponentially-weighted moving average:
\begin{equation}
  \bmu_c^T = (1-\eta)^T \bw_c + \eta \sum_{i=1}^T (1-\eta)^{T-i} w_i\, \bx_{t_i}
  \label{eq:drift}
\end{equation}
where $w_i = w_{\text{upd}}$ at step $i$.
The text anchor $\bw_c$ is the initialiser that decays geometrically; past samples accumulate with recency weighting.
The parameter $\eta$ controls how fast the anchor is forgotten.

\subsubsection{Logit Fusion}

The combined prediction is an adaptive fusion of two classifiers:
\begin{equation}
  \mathbf{z}^{\text{combined}} = \underbrace{\mathbf{z}^{\text{clip}}}_{\text{text prior}} + \alpha_t \cdot \underbrace{\mathbf{z}^{\text{adapt}}}_{\text{visual centroid}}
\end{equation}
where $\alpha_t = 1 - H_{\text{clip}}(\bx_t)/\log C \in [0,1]$.
When CLIP is confident ($\alpha_t \approx 1$), the adapted mean dominates.
When CLIP is uncertain ($\alpha_t \approx 0$), the text anchor dominates.

\subsection{Hyperparameters}

\begin{table}[H]
\centering
\caption{FreeTTA hyperparameters.}
\label{tab:freetta_hyperparams}
\begin{tabular}{@{}lllll@{}}
\toprule
Parameter & Symbol & Value & Meaning \\
\midrule
Update step size   & $\eta$          & 0.10  & Mean shift per M-step \\
Adapted temperature & $\tau_{\text{adapt}}$ & 100 & Temperature for adapted logits \\
Gate steepness     & $\gamma$        & 10    & Sigmoid slope for update weighting \\
\bottomrule
\end{tabular}
\end{table}

\subsection{Key Structural Properties}

\begin{itemize}
  \item \textbf{Global}: predictions are informed by \emph{all} past samples for the predicted class.
  \item \textbf{No cold-start}: CLIP text anchor provides reasonable initialisation at step 0.
  \item \textbf{Uncertainty-aware}: low-confidence samples contribute small M-step updates ($w_{\text{upd}} \approx 0$).
  \item \textbf{Risk of drift}: if many early samples are misclassified, $\bmu_{\hat{y}}$ (wrong class) moves toward a different class's region, propagating errors.
\end{itemize}

% ══════════════════════════════════════════════════════════════════════════════
\section{How Much Do They Solve the Problem?}
% ══════════════════════════════════════════════════════════════════════════════

\subsection{Quantitative Summary}

\begin{table}[H]
\centering
\caption{Accuracy and error reduction across all datasets. Error reduction = $(a_\text{method} - a_\text{CLIP})/(1 - a_\text{CLIP}) \times 100\%$.}
\label{tab:accuracy_summary}
\begin{tabular}{@{}lrrrrrr@{}}
\toprule
Dataset & CLIP & TDA & FreeTTA & TDA Gain & FreeTTA Gain & Err.\ Reduction (FreeTTA) \\
\midrule
EuroSAT     & 48.43\% & 53.33\% & 59.35\% & +4.90\% & +10.91\% & 21.2\% \\
DTD         & 43.94\% & 45.85\% & 45.48\% & +1.91\% & +1.54\%  & 2.7\%  \\
Caltech101  & 93.55\% & 93.43\% & 93.63\% & $-$0.12\% & +0.08\% & 1.2\% \\
ImageNet-V2 & 62.35\% & 62.72\% & 62.72\% & +0.37\% & +0.37\%  & 0.98\% \\
Oxford Pets & 88.39\% & 88.42\% & 88.96\% & +0.03\% & +0.57\%  & 4.9\%  \\
\bottomrule
\end{tabular}
\end{table}

\subsection{Qualitative Assessment}

Both methods solve the problem only partially.
The key insight is: \textbf{TTA is most valuable exactly where CLIP fails most (high domain shift).}
On EuroSAT, FreeTTA corrects 21.2\% of CLIP's errors.
On Caltech, where CLIP already achieves 93.55\%, barely 1\% of errors are fixable.
This is self-reinforcing: the more domain shift, the more information the test stream carries about the true distribution.

% ══════════════════════════════════════════════════════════════════════════════
\section{Datasets Explained}
% ══════════════════════════════════════════════════════════════════════════════

\subsection{Caltech101}
101 object categories including animals, vehicles, and everyday objects in standard photography.
Caltech classes appear frequently in internet text--image pairs, so CLIP's training distribution closely matches the visual style.
With 93.55\% baseline, there is very little room to improve.
\textit{Key stats}: 2,465 samples, 100 classes ($\approx$24.7/class), cache pressure 4.93$\times$, GAS $=+0.056$.

\subsection{DTD (Describable Textures Dataset)}
47 texture categories (braided, bumpy, cobwebbed, cracked, etc.) photographed as close-up material surfaces.
Texture names are semantically similar in language, clustering CLIP anchors closely together.
High intra-class variance makes it hard for any method to achieve high accuracy.
\textit{Key stats}: 1,880 samples, 47 classes ($\approx$40/class), cache pressure 8.0$\times$, GAS $=+0.096$.

\subsection{EuroSAT}
Satellite imagery of 10 European land-use categories.
The most extreme domain shift: overhead multispectral tiles versus internet ground-level photos.
CLIP text anchors encode terrestrial perspectives; test images are $64\times64$ aerial tiles.
\textit{Key stats}: 8,100 samples, 10 classes (810/class), cache pressure 162$\times$, GAS $=-0.115$.

\subsection{ImageNet-V2}
10,000 images from 1,000 ImageNet classes (10 per class), selected as a harder test set.
With 1,000 classes and only 10 samples each, neither method can accumulate meaningful class statistics.
\textit{Key stats}: cache pressure 2.0$\times$, GAS $=+0.408$ (centroid oracle 82.7\% vs 1-NN oracle 41.9\%).

\subsection{Oxford Pets}
37 dog and cat breed categories.
Challenge: fine-grained breed discrimination (Samoyed vs American Eskimo Dog vs Bichon Frise).
\textit{Key stats}: 3,669 samples, 37 classes ($\approx$99/class), cache pressure 19.8$\times$, GAS $=+0.080$.

% ══════════════════════════════════════════════════════════════════════════════
\section{Comparison Approach and Metrics}
% ══════════════════════════════════════════════════════════════════════════════

\subsection{Research Philosophy}

\textbf{Central question}: Given the same frozen CLIP features, does it matter more to remember individual examples (TDA) or to update a global class model (FreeTTA)?

Accuracy alone cannot answer this.
A 1\% difference could arise from many different mechanisms.
We designed metrics that trace \emph{how} each method changes CLIP's predictions, \emph{with what quality}, \emph{at what speed}, and \emph{in what failure modes}.

The eight diagnostic questions we asked:
\begin{enumerate}
  \item Does the method change anything at all? $\to$ Change Rate, Flip Analysis
  \item Are the changes good or bad? $\to$ Beneficial Flip Precision (BFP)
  \item When does it start helping? $\to$ Break-Even Point
  \item How certain is the method? $\to$ Entropy \& Confidence
  \item When do the two methods disagree? $\to$ Disagreement Analysis
  \item What types of failures remain? $\to$ Failure Buckets
  \item Is the feature space geometry suited to each method? $\to$ Geometry Alignment Score (GAS)
  \item How hard is the adaptation problem? $\to$ Cache Pressure Ratio, EM Weight
\end{enumerate}

\subsection{Metric Definitions}

\subsubsection{Metric 1: Accuracy}
\begin{equation}
  \text{Accuracy} = \frac{1}{N}\sum_{t=1}^N \mathbf{1}\!\left[\hat{y}_t = y_t\right]
\end{equation}

The fundamental outcome measure. Tells us the headline but not the mechanism.

\subsubsection{Metric 2: Change Rate and Flip Analysis}

\begin{equation}
  \text{Change Rate} = \frac{1}{N}\sum_{t=1}^N \mathbf{1}\!\left[\hat{y}_t^{\text{method}} \neq \hat{y}_t^{\text{clip}}\right]
\end{equation}

Each changed prediction is categorised:
\begin{align}
  \text{Beneficial flip:} &\quad \hat{y}^{\text{clip}} \neq y \;\text{ AND }\; \hat{y}^{\text{method}} = y \\
  \text{Harmful flip:} &\quad \hat{y}^{\text{clip}} = y \;\text{ AND }\; \hat{y}^{\text{method}} \neq y \\
  \text{Other changed wrong:} &\quad \hat{y}^{\text{clip}} \neq y \;\text{ AND }\; \hat{y}^{\text{method}} \neq y \;\text{ AND }\; \hat{y}^{\text{method}} \neq \hat{y}^{\text{clip}}
\end{align}
\begin{equation}
  \text{Net Correction Rate} = \frac{N_{\text{ben}} - N_{\text{harm}}}{N}
\end{equation}

\subsubsection{Metric 3: Beneficial Flip Precision (BFP)}
\begin{equation}
  \text{BFP} = \frac{N_{\text{beneficial}}}{N_{\text{beneficial}} + N_{\text{harmful}} + N_{\text{other changed wrong}}}
             = \frac{N_{\text{beneficial}}}{N_{\text{total changed}}}
  \label{eq:bfp}
\end{equation}

Fraction of all prediction changes that are improvements.
BFP $= 1.0$: every change is beneficial.
BFP $= 0.5$: changes are random noise.
BFP $< 0.33$: method is actively harmful on average.

\subsubsection{Metric 4: Entropy and Confidence}
\begin{equation}
  H(\bx) = -\sum_{c=1}^C p(c \mid \bx)\,\log p(c \mid \bx)
  \label{eq:entropy}
\end{equation}
\begin{equation}
  h(\bx) = \frac{H(\bx)}{\log C} \in [0, 1] \quad\text{(normalised entropy)}
\end{equation}
\begin{equation}
  \text{Confidence}(\bx) = \max_c p(c \mid \bx)
\end{equation}

We track entropy and confidence separately for correct and wrong predictions.
Well-calibrated methods have low entropy on correct predictions and high entropy on wrong ones.

\subsubsection{Metric 5: Break-Even Point}
\begin{equation}
  \text{Break-Even} = \min\left\{t \;:\; \overline{\text{Acc}}_{\text{method}}(t) > \overline{\text{Acc}}_{\text{clip}}(t)\right\}
\end{equation}
where $\overline{\text{Acc}}(t)$ is the rolling accuracy over the last $W=50$ steps.
A lower break-even ratio (step / total steps) means the method helps sooner.

\subsubsection{Metric 6: Disagreement Analysis}
\begin{equation}
  \text{Disagreement Rate} = \frac{1}{N}\sum_t \mathbf{1}\!\left[\hat{y}_t^{\text{TDA}} \neq \hat{y}_t^{\text{FreeTTA}}\right]
\end{equation}
\begin{equation}
  \text{Acc}_{\text{method on disagree}} = \frac{1}{|\mathcal{D}|}\sum_{t \in \mathcal{D}} \mathbf{1}\!\left[\hat{y}_t^{\text{method}} = y_t\right]
\end{equation}
where $\mathcal{D} = \{t : \hat{y}_t^{\text{TDA}} \neq \hat{y}_t^{\text{FreeTTA}}\}$.

\subsubsection{Metric 7: Failure Buckets}

Every sample falls into exactly one of:

\begin{table}[H]
\centering
\caption{Failure bucket definitions.}
\label{tab:buckets_def}
\begin{tabular}{@{}ll@{}}
\toprule
Bucket & Meaning \\
\midrule
\texttt{only\_freetta\_correct}    & $\hat{y}^{\text{clip}} \neq y$, $\hat{y}^{\text{tda}} \neq y$, $\hat{y}^{\text{ftta}} = y$ \\
\texttt{only\_tda\_correct}        & $\hat{y}^{\text{clip}} \neq y$, $\hat{y}^{\text{ftta}} \neq y$, $\hat{y}^{\text{tda}} = y$ \\
\texttt{clip\_correct\_both\_fail} & $\hat{y}^{\text{clip}} = y$, $\hat{y}^{\text{tda}} \neq y$, $\hat{y}^{\text{ftta}} \neq y$ \\
\texttt{all\_wrong}                & $\hat{y}^{\text{clip}} \neq y$, $\hat{y}^{\text{tda}} \neq y$, $\hat{y}^{\text{ftta}} \neq y$ \\
\bottomrule
\end{tabular}
\end{table}

\subsubsection{Metric 8: Geometry Alignment Score (GAS)}
\begin{equation}
  \text{GAS} = \underbrace{\text{OracleCentroidAcc}}_{\text{upper bound for FreeTTA}} - \underbrace{\text{Oracle1NNAcc}}_{\text{upper bound for TDA}}
  \label{eq:gas}
\end{equation}
\begin{itemize}[noitemsep]
  \item $\text{GAS} > 0$: true class centroids are better discriminators than individual instances $\to$ FreeTTA favoured.
  \item $\text{GAS} < 0$: local neighbourhood structure is stronger than global means $\to$ TDA favoured.
\end{itemize}

\subsubsection{Metric 9: Cache Pressure and EM Weight}
\begin{equation}
  \text{Cache Pressure} = \frac{N_{\text{total}}}{C \times K^+}
  \label{eq:cache_pressure}
\end{equation}
\begin{equation}
  \overline{\alpha} = \frac{1}{N}\sum_{t=1}^N \alpha_t = \frac{1}{N}\sum_{t=1}^N \left(1 - \frac{H_{\text{clip}}(\bx_t)}{\log C}\right)
  \label{eq:em_weight}
\end{equation}
High cache pressure predicts TDA failure.
Low $\overline{\alpha}$ predicts slow FreeTTA adaptation.

% ══════════════════════════════════════════════════════════════════════════════
\section{Results: Dataset-by-Dataset}
% ══════════════════════════════════════════════════════════════════════════════

\subsection{Caltech101}

\begin{table}[H]
\centering
\caption{Caltech101 --- accuracy and flip metrics.}
\label{tab:caltech_flip}
\begin{tabular}{@{}lrrrrrr@{}}
\toprule
Method & Acc.\ & Gain & Change Rate & Ben.\ Flips & Harm.\ Flips & BFP \\
\midrule
CLIP    & 93.55\% & ---    & 0.00\% & --- & --- & --- \\
TDA     & 93.43\% & $-$0.12\% & 1.14\% & 12  & 15  & 42.9\% \\
FreeTTA & 93.63\% & +0.08\% & 0.41\% & 6   & 4   & 60.0\% \\
\bottomrule
\end{tabular}
\end{table}

\begin{table}[H]
\centering
\caption{Caltech101 --- entropy and confidence.}
\label{tab:caltech_entropy}
\begin{tabular}{@{}lrrrr@{}}
\toprule
Method & Correct conf & Wrong conf & Correct entropy & Wrong entropy \\
\midrule
CLIP    & 0.0113 & 0.0112 & 4.605 & 4.605 \\
TDA     & 0.9690 & 0.7476 & 0.130 & 0.724 \\
FreeTTA & 0.9271 & 0.6168 & 0.293 & 1.071 \\
\bottomrule
\end{tabular}
\end{table}

\begin{table}[H]
\centering
\caption{Caltech101 --- key internal metrics and failure buckets.}
\label{tab:caltech_internal}
\begin{tabular}{@{}lr@{}}
\toprule
Metric & Value \\
\midrule
Cache pressure ratio        & 4.93 \\
Mean EM weight $\overline{\alpha}$ & 0.826 \\
FreeTTA final $\mu$ drift   & 1.134 \\
GAS                         & $+0.056$ \\
Oracle centroid acc         & 97.44\% \\
Oracle 1-NN acc             & 91.85\% \\
Disagreement rate           & 1.14\% \\
TDA acc on disagreements    & 39.3\% \\
FreeTTA acc on disagreements & 57.1\% \\
Break-even TDA              & N/A (TDA ends below CLIP) \\
Break-even FreeTTA          & Step 1050 / 2465 (42.6\%) \\
\midrule
\texttt{only\_freetta\_correct} & 3 (0.12\%) \\
\texttt{only\_tda\_correct}     & 9 (0.37\%) \\
\texttt{clip\_correct, tda\_wrong} & 13 (0.53\%) \\
\texttt{clip\_correct, ftta\_wrong} & 2 (0.08\%) \\
\texttt{all\_wrong}             & 144 (5.84\%) \\
\bottomrule
\end{tabular}
\end{table}

\textbf{Analysis.} Caltech is near-solved by CLIP.
Both adapters make very few changes.
FreeTTA is more conservative (0.41\% vs 1.14\% change rate) and more precise (BFP~60\% vs 43\%).
TDA makes more changes but also more harmful ones---15 harmful vs FreeTTA's 4---ending slightly below CLIP ($-$0.12\%).
GAS $= +0.056$ correctly predicts FreeTTA edges out TDA (global centroids slightly better than local lookups).
FreeTTA breaks even at step 1050 (42.6\% into the stream); despite this late break-even, FreeTTA finishes with a net gain (+0.08\%), while TDA never reaches a net benefit over CLIP.

\subsection{DTD (Describable Textures)}

\begin{table}[H]
\centering
\caption{DTD --- accuracy and flip metrics.}
\label{tab:dtd_flip}
\begin{tabular}{@{}lrrrrrr@{}}
\toprule
Method & Acc.\ & Gain & Change Rate & Ben.\ Flips & Harm.\ Flips & BFP \\
\midrule
CLIP    & 43.94\% & ---    & 0.00\%  & ---  & ---  & --- \\
TDA     & 45.85\% & +1.91\% & 18.51\% & 87  & 51   & 25.0\% \\
FreeTTA & 45.48\% & +1.54\% & 8.09\%  & 44  & 15   & 28.9\% \\
\bottomrule
\end{tabular}
\end{table}

\begin{table}[H]
\centering
\caption{DTD --- key metrics.}
\label{tab:dtd_internal}
\begin{tabular}{@{}lr@{}}
\toprule
Metric & Value \\
\midrule
Cache pressure ratio         & 8.0 \\
Mean EM weight               & 0.253 \\
GAS                          & $+0.096$ \\
Oracle centroid / 1-NN acc   & 73.19\% / 63.62\% \\
Disagreement rate            & 17.23\% \\
TDA / FreeTTA acc on disagree & 21.3\% / 19.1\% \\
Break-even TDA / FreeTTA     & Step 118 (6.3\%) / Step 9 (0.5\%) \\
\texttt{only\_freetta\_correct} & 22 (1.17\%) \\
\texttt{only\_tda\_correct}     & 65 (3.46\%) \\
\texttt{clip\_correct, tda\_wrong}  & 51 (2.71\%) \\
\texttt{clip\_correct, ftta\_wrong} & 15 (0.80\%) \\
\texttt{all\_wrong}             & 945 (50.27\%) \\
\bottomrule
\end{tabular}
\end{table}

\textbf{Analysis.} DTD reveals the texture classification ceiling starkly: $\approx$50\% of all samples are irreducible errors given frozen CLIP features.
\textbf{TDA outperforms FreeTTA on DTD} (+1.91\% vs +1.54\%).
With a tighter entropy gate ($\theta_{\text{low}}=0.05$, $\theta_{\text{high}}=0.40$ vs the defaults of 0.20/0.50) and a larger positive cache ($K=5$ vs $K=3$), TDA more selectively admits high-confidence texture exemplars and retains a richer set of per-class features.
The discrete, non-overlapping nature of texture categories (``bumpy'' vs ``dotted'' vs ``woven'') means that confident local cached examples transfer reliably to future test samples from the same class.
FreeTTA's global class mean is pulled in conflicting directions by the high visual ambiguity among texture categories, limiting its adapted prototypes.
GAS $= +0.096$ nominally favours centroid-based methods, yet TDA's precise entropy filtering produces more accurate local exemplars than FreeTTA's weighted mean update.

\subsection{EuroSAT}

\begin{table}[H]
\centering
\caption{EuroSAT --- accuracy and flip metrics.}
\label{tab:eurosat_flip}
\begin{tabular}{@{}lrrrrrr@{}}
\toprule
Method & Acc.\ & Gain & Change Rate & Ben.\ Flips & Harm.\ Flips & BFP \\
\midrule
CLIP    & 48.43\% & ---     & 0.00\%  & ---   & ---  & --- \\
TDA     & 53.33\% & +4.90\% & 41.91\% & 1,125 & 728  & 33.1\% \\
FreeTTA & 59.35\% & +10.91\% & 30.95\% & 1,258 & 374  & 50.2\% \\
\bottomrule
\end{tabular}
\end{table}

\begin{table}[H]
\centering
\caption{EuroSAT --- entropy and confidence.}
\label{tab:eurosat_entropy}
\begin{tabular}{@{}lrrrr@{}}
\toprule
Method & Correct conf & Wrong conf & Correct entropy & Wrong entropy \\
\midrule
CLIP    & 0.1036 & 0.1020 & 2.302 & 2.302 \\
TDA     & 0.7320 & 0.5600 & 0.776 & 1.207 \\
FreeTTA & 0.8119 & 0.5195 & 0.512 & 1.227 \\
\bottomrule
\end{tabular}
\end{table}

\begin{table}[H]
\centering
\caption{EuroSAT --- key metrics.}
\label{tab:eurosat_internal}
\begin{tabular}{@{}lr@{}}
\toprule
Metric & Value \\
\midrule
Cache pressure ratio         & 162.0 \\
Mean EM weight               & 0.198 \\
GAS                          & $-0.115$ \\
Oracle centroid / 1-NN acc   & 78.26\% / 89.80\% \\
Disagreement rate            & 30.90\% \\
TDA / FreeTTA acc on disagree & 22.8\% / 42.3\% \\
Break-even TDA / FreeTTA     & Step 941 (11.6\%) / Step 8 (0.10\%) \\
\texttt{only\_freetta\_correct} & 621 (7.67\%) \\
\texttt{only\_tda\_correct}     & 488 (6.02\%) \\
\texttt{clip\_correct, tda\_wrong} & 83 (1.02\%) \\
\texttt{clip\_correct, ftta\_wrong} & 437 (5.40\%) \\
\texttt{all\_wrong}             & 2,431 (30.01\%) \\
\bottomrule
\end{tabular}
\end{table}

\textbf{Analysis.}
EuroSAT is the most revealing dataset.
\textbf{FreeTTA's +10.91\% is the largest absolute gain of any method on any dataset.}
Cache pressure of 162$\times$ renders TDA's memory nearly useless---each cache slot is overwritten 162 times per class.
TDA's 728 harmful flips (breaking CLIP's correct predictions) cost it dearly.
FreeTTA's BFP $= 50.2\%$ means every change it makes is a net positive on average.

The GAS $= -0.115$ is the only negative value across all datasets: local neighbourhood structure is better than global centroids on EuroSAT (oracle 1-NN $= 89.8\%$ vs centroid $= 78.3\%$).
Yet FreeTTA still wins because the \emph{initial} CLIP text anchor is so far from the true visual distribution that any mean correction outweighs the loss from using centroids instead of local lookups.

Mean EM weight $= 0.198$ (lowest across all datasets): CLIP is very uncertain on EuroSAT, so FreeTTA adapts cautiously---yet it still dominates by 6.02 points over TDA.

\subsection{ImageNet-V2}

\begin{table}[H]
\centering
\caption{ImageNet-V2 --- accuracy and flip metrics.}
\label{tab:imagenet_flip}
\begin{tabular}{@{}lrrrrrr@{}}
\toprule
Method & Acc.\ & Gain & Change Rate & Ben.\ Flips & Harm.\ Flips & BFP \\
\midrule
CLIP    & 62.35\% & ---    & 0.00\% & ---  & ---  & --- \\
TDA     & 62.72\% & +0.37\% & 3.88\% & 115  & 78   & 29.6\% \\
FreeTTA & 62.72\% & +0.37\% & 3.62\% & 116  & 79   & 32.0\% \\
\bottomrule
\end{tabular}
\end{table}

\begin{table}[H]
\centering
\caption{ImageNet-V2 --- key metrics.}
\label{tab:imagenet_internal}
\begin{tabular}{@{}lr@{}}
\toprule
Metric & Value \\
\midrule
Cache pressure ratio         & 2.0 \\
Mean EM weight               & 0.533 \\
GAS                          & $+0.408$ \\
Oracle centroid / 1-NN acc   & 82.72\% / 41.89\% \\
Disagreement rate            & 5.60\% \\
TDA / FreeTTA acc on disagree & 25.2\% / 25.2\% \\
Break-even TDA / FreeTTA     & Step 363 (3.6\%) / Step 2 (0.02\%) \\
\texttt{only\_freetta\_correct} & 81 (0.81\%) \\
\texttt{only\_tda\_correct}     & 80 (0.80\%) \\
\texttt{all\_wrong}             & 3,569 (35.7\%) \\
\bottomrule
\end{tabular}
\end{table}

\textbf{Analysis.}
A statistical dead tie.
With only 10 samples per class across 1,000 classes, neither method accumulates enough information to adapt.
GAS $= +0.408$ is the largest positive value---centroid oracle 82.7\% vs 1-NN oracle 41.9\%---strongly predicting FreeTTA should dominate.
Yet the geometry is irrelevant without sufficient samples: FreeTTA's mean estimates never converge with 10 updates per class.
Identical 25.2\% accuracy on disagreements confirms both methods are effectively coin-flipping on hard cases.

\subsection{Oxford Pets}

\begin{table}[H]
\centering
\caption{Oxford Pets --- accuracy and flip metrics.}
\label{tab:pets_flip}
\begin{tabular}{@{}lrrrrrr@{}}
\toprule
Method & Acc.\ & Gain & Change Rate & Ben.\ Flips & Harm.\ Flips & BFP \\
\midrule
CLIP    & 88.39\% & ---     & 0.00\% & ---  & ---  & --- \\
TDA     & 88.42\% & +0.03\% & 3.35\% & 52   & 51   & 42.3\% \\
FreeTTA & 88.96\% & +0.57\% & 3.08\% & 56   & 35   & 49.6\% \\
\bottomrule
\end{tabular}
\end{table}

\begin{table}[H]
\centering
\caption{Oxford Pets --- key metrics.}
\label{tab:pets_internal}
\begin{tabular}{@{}lr@{}}
\toprule
Metric & Value \\
\midrule
Cache pressure ratio         & 19.8 \\
Mean EM weight               & 0.650 \\
GAS                          & $+0.080$ \\
Oracle centroid / 1-NN acc   & 91.50\% / 83.46\% \\
Disagreement rate            & 3.30\% \\
TDA / FreeTTA acc on disagree & 35.5\% / 52.1\% \\
Break-even TDA / FreeTTA     & Step 229 (6.2\%) / Step 5 (0.1\%) \\
\texttt{only\_freetta\_correct} & 32 (0.87\%) \\
\texttt{only\_tda\_correct}     & 28 (0.76\%) \\
\texttt{clip\_correct, tda\_wrong} & 31 (0.84\%) \\
\texttt{clip\_correct, ftta\_wrong} & 15 (0.41\%) \\
\texttt{all\_wrong}             & 342 (9.32\%) \\
\bottomrule
\end{tabular}
\end{table}

\textbf{Analysis.}
\textbf{FreeTTA outperforms TDA on Pets} (+0.57\% vs +0.03\%), with TDA barely clearing CLIP.
With $\alpha=0.25$ and $\beta=4.0$, FreeTTA's adapted class means capture the fine-grained breed distribution more effectively than TDA's local cache.
GAS $= +0.080$ correctly predicts the outcome: oracle centroids (91.50\%) outperform oracle 1-NN (83.46\%), indicating that global breed prototypes are more discriminative than local nearest-neighbor lookups for the 37 Oxford Pets breeds.
FreeTTA's higher precision (BFP~49.6\% vs 42.3\%) reflects that its global mean adaptation corrects more CLIP errors: 32 samples CLIP missed but FreeTTA fixed (vs.\ 28 for TDA), with 15 harmful CLIP-correct flips (vs.\ 51 for TDA).
TDA's exact cached features provide negligible benefit (+0.03\% over CLIP): with only 3 cache slots per class and $\sim$99 test images per breed, cache pressure is high (ratio 19.8), leading to near-equal beneficial and harmful flips (52 vs.\ 51).

% ══════════════════════════════════════════════════════════════════════════════
\section{Cross-Dataset Conclusions}
% ══════════════════════════════════════════════════════════════════════════════

\subsection{Complete Metrics Summary}

\begin{table}[H]
\centering
\caption{All datasets: complete metric summary.}
\label{tab:full_summary}
\setlength{\tabcolsep}{4pt}
\begin{tabular}{@{}lrrrrrrrrrl@{}}
\toprule
Dataset & CLIP & TDA & FreeTTA & BFP$_\text{TDA}$ & BFP$_\text{Ftta}$ & GAS & CacheP & $\overline{\alpha}$ & Disagree & Ftta wins? \\
\midrule
Caltech  & 93.55 & 93.43 & 93.63 & 42.9\% & 60.0\% & $+0.056$ & 4.93  & 0.826 & 1.14\%  & Yes \\
DTD      & 43.94 & 45.85 & 45.48 & 25.0\% & 28.9\% & $+0.096$ & 8.0   & 0.253 & 17.23\% & No \\
EuroSAT  & 48.43 & 53.33 & 59.35 & 33.1\% & 50.2\% & $-0.115$ & 162.0 & 0.198 & 30.90\% & Yes \\
ImageNet & 62.35 & 62.72 & 62.72 & 29.6\% & 32.0\% & $+0.408$ & 2.0   & 0.533 & 5.60\%  & Tie \\
Pets     & 88.39 & 88.42 & 88.96 & 42.3\% & 49.6\% & $+0.080$ & 19.8  & 0.650 & 3.30\%  & Yes \\
\bottomrule
\end{tabular}
\end{table}

\subsection{Key Findings}

\begin{tcolorbox}[findingbox, title={Finding 1: FreeTTA is the better default method.}]
FreeTTA wins or ties on 4 of 5 datasets (Caltech, EuroSAT, ImageNet, Pets).
Its advantage is largest where it matters most: EuroSAT +10.91\%.
Its only loss is DTD, where TDA's tighter entropy filtering (+1.91\% vs +1.54\%) better captures discrete texture categories.
\end{tcolorbox}

\begin{tcolorbox}[findingbox, title={Finding 2: Cache pressure correlates with FreeTTA's advantage (with one exception).}]
\begin{center}
\begin{tabular}{lrc}
\toprule
Dataset & Cache Pressure & FreeTTA $-$ TDA gap \\
\midrule
ImageNet & 2.0   & +0.00\% (tie) \\
Caltech  & 4.93  & +0.20\% \\
DTD      & 8.0   & $\mathbf{-0.37\%}$ \textbf{(TDA wins)} \\
Pets     & 19.8  & +0.54\% \\
EuroSAT  & 162.0 & +6.01\% \\
\bottomrule
\end{tabular}
\end{center}
Higher cache pressure generally predicts FreeTTA's advantage, but DTD is an exception: the discrete, non-overlapping nature of texture categories allows TDA's cache to transfer reliably despite moderate pressure.
\end{tcolorbox}

\begin{tcolorbox}[findingbox, title={Finding 3: BFP consistently favours FreeTTA.}]
On all 5 datasets, FreeTTA's BFP $>$ TDA's BFP (60.0\% vs 42.9\%, 28.9\% vs 25.0\%, 50.2\% vs 33.1\%, 32.0\% vs 29.6\%, 49.6\% vs 42.3\%).
FreeTTA's changes are more often improvements---its adaptation signal is consistently higher quality.
On DTD, despite TDA's lower BFP, TDA wins on accuracy because it makes far more changes in absolute terms (18.51\% vs 8.09\% change rate), yielding 87 vs 44 beneficial flips.
\end{tcolorbox}

\begin{tcolorbox}[findingbox, title={Finding 4: FreeTTA adapts faster on all datasets.}]
\begin{center}
\begin{tabular}{lrr}
\toprule
Dataset & TDA Break-Even & FreeTTA Break-Even \\
\midrule
EuroSAT  & Step 941  (11.6\%) & Step 8   (0.10\%) \\
DTD      & Step 118  (6.3\%)  & Step 9   (0.50\%) \\
Caltech  & N/A (TDA $<$ CLIP) & Step 1050 (42.6\%) \\
ImageNet & Step 363  (3.6\%)  & Step 2   (0.02\%) \\
Pets     & Step 229  (6.2\%)  & Step 5   (0.14\%) \\
\bottomrule
\end{tabular}
\end{center}
FreeTTA breaks even earlier on every dataset including DTD (step 9 vs 118).
TDA's accuracy win on DTD is not due to faster adaptation but to higher correction volume:
TDA makes 18.5\% change rate (87 beneficial flips) vs FreeTTA's 8.1\% (44 beneficial flips).
On Caltech, TDA never reaches break-even (ends below CLIP).
\end{tcolorbox}

\begin{tcolorbox}[findingbox, title={Finding 5: Both methods fix CLIP's broken calibration.}]
CLIP: correct confidence 0.0113, wrong confidence 0.0112 (indistinguishable!).\\
After TDA: 0.969 (correct) vs 0.748 (wrong). After FreeTTA: 0.927 vs 0.617.\\
Both dramatically improve calibration.
FreeTTA's wrong-prediction confidence is slightly higher---a signature of mean drift.
\end{tcolorbox}

% ══════════════════════════════════════════════════════════════════════════════
\section{Deep Analysis}
% ══════════════════════════════════════════════════════════════════════════════

% ── 1. Adaptation Dynamics ──────────────────────────────────────────────────
\subsection{Adaptation Dynamics}

\subsubsection{Rolling Accuracy Over Stream}

Rather than just comparing final accuracy, we track \emph{how} each method improves over the stream
using a rolling window of $w=50$ samples.
Table~\ref{tab:break_even_rolling} reports the first step at which the rolling accuracy of each
adapter surpasses CLIP's rolling accuracy in the same window.

\begin{table}[H]
\centering
\caption{Rolling break-even point (window $= 50$): first step where the adapter's rolling accuracy
exceeds CLIP's rolling accuracy.}
\label{tab:break_even_rolling}
\begin{tabular}{@{}lrrrr@{}}
\toprule
Dataset & $N$ & TDA break-even & FreeTTA break-even & FreeTTA advantage \\
\midrule
Caltech  & 2,465  & N/A (ends below CLIP) & Step 1050 (42.6\%)  & FreeTTA $>$ CLIP; TDA $<$ CLIP \\
DTD      & 1,880  & Step 118 (6.3\%)       & Step 9 (0.5\%)      & $5.7\%$ faster \\
EuroSAT  & 8,100  & Step 941 (11.6\%)      & Step 8 (0.1\%)      & $11.5\%$ faster \\
Pets     & 3,669  & Step 229 (6.2\%)       & Step 5 (0.1\%)      & $6.1\%$ faster \\
ImageNet & 10,000 & Step 363 (3.6\%)       & Step 2 (0.02\%)     & $3.6\%$ faster \\
\bottomrule
\end{tabular}
\end{table}

FreeTTA breaks even within the first $0.5\%$ of the stream on every dataset except Caltech.
TDA requires up to $11.6\%$ of the stream before its rolling accuracy equals CLIP's---a \textbf{cold-start penalty} caused by the empty cache.
On Caltech, TDA's rolling accuracy never consistently exceeds CLIP's: with only 15 harmful flips vs 12 beneficial across the full stream, TDA ends slightly below CLIP.

\subsubsection{Adaptation Speed: Mechanisms}

\textbf{Why FreeTTA adapts instantly.}
At step $t=1$ the adapted mean $\bmu_c$ has already shifted by
$\eta \cdot w_\text{upd}(\bx_1 - \bmu_c^{(0)})$.
Even this single-sample update is enough to improve over the frozen CLIP anchor if
$\bx_1$ is confidently classified (large $w_\text{upd}$), because $\bmu_c^{(0)}$ starts
at the text anchor, which is already a reasonable initialisation.

\textbf{Why TDA suffers cold-start.}
TDA's positive cache for class $c$ is empty at $t=1$.
With no cached features, the adjustment $z^\text{pos}_c = 0$, so TDA reduces to CLIP for all
classes not yet encountered in the stream.
Once the cache fills, TDA's corrections become reliable---but filling requires $C \cdot K_\text{pos}$ high-confidence samples (e.g.\ $47 \times 5 = 235$ on DTD).

\subsubsection{Stability Over Time}

Table~\ref{tab:stability} reports the standard deviation of rolling accuracy in the last 30\% of
each stream (see \texttt{outputs/dynamics\_analysis/adaptation\_rolling\_accuracy.png} for curves).

\begin{table}[H]
\centering
\caption{Late-stream rolling-accuracy stability: $\sigma$ of rolling acc (\%) in final 30\% of stream.
Lower = more stable.}
\label{tab:stability}
\begin{tabular}{@{}lrrrr@{}}
\toprule
Dataset & CLIP $\sigma$ & TDA $\sigma$ & FreeTTA $\sigma$ & Best \\
\midrule
Caltech  & 8.8\%  & 8.6\%  & 8.9\%  & TDA $\approx$ FreeTTA \\
DTD      & 16.2\% & 18.0\% & 16.7\% & CLIP (neither helps) \\
EuroSAT  & 34.9\% & 19.7\% & 17.8\% & FreeTTA (major stabilisation) \\
Pets     & 7.2\%  & 9.4\%  & 7.3\%  & CLIP $\approx$ FreeTTA \\
ImageNet & 11.5\% & 11.2\% & 11.3\% & Tie \\
\bottomrule
\end{tabular}
\end{table}

Key observations:
\begin{itemize}
  \item \textbf{EuroSAT is the stand-out case}: both adapters reduce rolling-accuracy variance from
  34.9\% (CLIP) to $\sim$18--20\%.
  Domain shift is so severe that CLIP's predictions fluctuate wildly between windows;
  both adapters learn a stable domain representation that smooths these oscillations.
  FreeTTA achieves the lowest variance (17.8\%), because its global mean is updated across all 810
  samples per class.
  \item \textbf{DTD} shows slightly higher TDA variance (18.0\% vs 16.7\% FreeTTA): TDA's cache
  corrections occasionally mis-fire on new texture classes before they accumulate sufficient exemplars.
  \item On near-solved datasets (Caltech, Pets, ImageNet), stability is similar across all three
  methods---the adapters neither help nor hurt.
\end{itemize}

\subsubsection{FreeTTA Class-Mean ($\mu$) Drift}

Table~\ref{tab:mu_drift} shows the final Euclidean drift of adapted class means from their text
anchor initialisations, alongside the mean EM weight $\bar{\alpha}$.

\begin{table}[H]
\centering
\caption{FreeTTA internal adaptation state at end of stream.
$\bar{\alpha}$ = mean EM update weight (CLIP confidence proxy).
$\|\bmu - \bmu_0\|$ = final mean drift from text-anchor initialisation.}
\label{tab:mu_drift}
\begin{tabular}{@{}lrrrrr@{}}
\toprule
Dataset & $\bar{\alpha}$ & Final $\|\bmu - \bmu_0\|$ & $\sigma$-trace & Prior $H(\pi)$ & N/C \\
\midrule
Caltech  & 0.826 & 1.134 & 0.357 & 3.677 & 24.7 \\
DTD      & 0.253 & 1.155 & 0.333 & 1.251 & 40.0 \\
EuroSAT  & 0.198 & 1.142 & 0.348 & 0.717 & 810  \\
Pets     & 0.650 & 1.122 & 0.370 & 2.588 & 99.2 \\
ImageNet & 0.533 & 1.128 & 0.364 & 3.965 & 10   \\
\bottomrule
\end{tabular}
\end{table}

$\bar{\alpha}$ reflects how much useful signal the stream provides:
Caltech ($\bar{\alpha}=0.826$) is easy for CLIP, so each sample is high-confidence and M-steps are
large; EuroSAT ($\bar{\alpha}=0.198$) is hard for CLIP, so M-steps are small but persistent over
810 samples/class.
Remarkably, final $\mu$-drift values are similar ($1.12$--$1.16$) across all datasets, suggesting
the adapted means always travel roughly the same normalised distance from the text anchor,
regardless of how quickly they get there.

% ── 2. Uncertainty Analysis ──────────────────────────────────────────────────
\subsection{Uncertainty Analysis}

\subsubsection{CLIP's Calibration Failure}

CLIP produces \emph{near-uniformly distributed} softmax outputs on all five datasets (Table~\ref{tab:entropy_calibration}).
Its mean entropy for correct predictions and wrong predictions is virtually identical ($\Delta H \approx 0$),
making CLIP's raw confidence unreliable as a self-calibration signal for most datasets.

This is a structural consequence of the cosine-similarity scoring with fixed temperature $\tau=100$:
when CLIP text anchors are poorly aligned with test features, all logits are near-zero and softmax spreads probability uniformly.
On Caltech and Pets, CLIP is near-maximally confident but equally so for correct and wrong predictions.
On EuroSAT and DTD, even the mean entropy gap is essentially zero despite CLIP showing moderate rank-calibration
in Spearman terms (see Table~\ref{tab:spearman_rho})---the mean difference vanishes while within-distribution ordering remains.

\begin{table}[H]
\centering
\caption{Mean prediction entropy ($H$) and confidence for correct vs wrong predictions.
CLIP's near-zero entropy gap ($\Delta H$) confirms it cannot self-calibrate.}
\label{tab:entropy_calibration}
\begin{tabular}{@{}llrrrr@{}}
\toprule
Dataset & Method & $H_\text{correct}$ & $H_\text{wrong}$ & $\Delta H$ & Conf$_\text{correct}$ \\
\midrule
\multirow{3}{*}{Caltech}
 & CLIP    & 4.605 & 4.605 & 0.000 & 1.13\% \\
 & TDA     & 0.130 & 0.724 & 0.594 & 96.9\% \\
 & FreeTTA & 0.293 & 1.071 & 0.778 & 92.7\% \\
\midrule
\multirow{3}{*}{DTD}
 & CLIP    & 3.850 & 3.850 & 0.000 & 2.25\% \\
 & TDA     & 0.506 & 1.150 & 0.644 & 86.1\% \\
 & FreeTTA & 1.407 & 2.305 & 0.898 & 64.4\% \\
\midrule
\multirow{3}{*}{EuroSAT}
 & CLIP    & 2.302 & 2.302 & 0.000 & 10.4\% \\
 & TDA     & 0.776 & 1.207 & 0.431 & 73.2\% \\
 & FreeTTA & 0.512 & 1.227 & 0.715 & 81.2\% \\
\midrule
\multirow{3}{*}{Pets}
 & CLIP    & 3.610 & 3.610 & 0.000 & 3.05\% \\
 & TDA     & 0.289 & 0.939 & 0.650 & 91.1\% \\
 & FreeTTA & 0.302 & 1.002 & 0.700 & 90.0\% \\
\midrule
\multirow{3}{*}{ImageNet}
 & CLIP    & 6.907 & 6.907 & 0.000 & 0.12\% \\
 & TDA     & 0.733 & 1.792 & 1.059 & 83.0\% \\
 & FreeTTA & 0.850 & 1.861 & 1.011 & 77.8\% \\
\bottomrule
\end{tabular}
\end{table}

\subsubsection{Entropy Gap as Calibration Quality Metric}

Both adapters create a substantial entropy gap between correct and wrong predictions (all $\Delta H > 0.4$).
Spearman rank correlation between entropy and error indicator (Table~\ref{tab:spearman_rho}) confirms this:

\begin{table}[H]
\centering
\caption{Spearman $\rho$ (entropy vs.\ error indicator). Values near 0 = uncalibrated; higher = better calibration.}
\label{tab:spearman_rho}
\begin{tabular}{@{}lrrr@{}}
\toprule
Dataset & CLIP $\rho$ & TDA $\rho$ & FreeTTA $\rho$ \\
\midrule
Caltech  & $-0.076$ & $+0.360$ & $+0.332$ \\
DTD      & $+0.261$ & $+0.416$ & $+0.445$ \\
EuroSAT  & $+0.433$ & $+0.365$ & $+0.593$ \\
Pets     & $-0.028$ & $+0.419$ & $+0.435$ \\
ImageNet & $+0.074$ & $+0.499$ & $+0.493$ \\
\bottomrule
\end{tabular}
\end{table}

All adapter $\rho$ values have $p \ll 0.001$.
CLIP's Spearman $\rho$ is close to zero on Caltech ($-0.076$), Pets ($-0.028$), and ImageNet ($+0.074$)---consistent with its broken calibration (near-uniform entropy, $\Delta H \approx 0$).
Interestingly, CLIP retains moderate rank-calibration on DTD ($\rho = 0.261$) and EuroSAT ($\rho = 0.433$): despite the near-zero \emph{mean} entropy gap, there is systematic monotonic ordering within the distribution on hard datasets.
Both adapters substantially improve calibration; however, on EuroSAT, TDA's $\rho = 0.365$ is \emph{lower} than CLIP's $\rho = 0.433$---consistent with TDA's stale, high-pressure cache producing noisier logit corrections that actually hurt rank-ordering.
FreeTTA on EuroSAT achieves the highest $\rho = 0.593$ of any method on any dataset.
This is the core mechanism that makes downstream use of adapter confidence scores meaningful.
Key observations:
\begin{itemize}
  \item \textbf{FreeTTA has larger $\Delta H$ than TDA on 4 of 5 datasets}
  (Caltech: 0.778 vs 0.594; DTD: 0.898 vs 0.644; EuroSAT: 0.715 vs 0.431; Pets: 0.700 vs 0.650).
  This means FreeTTA's wrong predictions carry \emph{more} uncertainty, making them easier to detect.
  The exception is ImageNet, where TDA's $\Delta H = 1.059 > 1.011$ for FreeTTA---both adapters
  perform identically and TDA's cache gives slightly sharper predictions.

  \item \textbf{TDA achieves lower absolute entropy on correct predictions} on most datasets
  (Caltech: 0.130 vs 0.293; DTD: 0.506 vs 1.407; Pets: 0.289 vs 0.302).
  This reflects TDA's hard-lookup mechanism: once a matching cached feature is found, it produces
  a near-maximum-confidence logit, collapsing entropy.
  FreeTTA's soft mean update yields smoother, more calibrated distributions.

  \item \textbf{EuroSAT is the exception}: FreeTTA's correct entropy (0.512) is \emph{lower} than
  TDA's (0.776), because FreeTTA successfully adapts its means and becomes more confident.
  TDA's cache is under 162$\times$ pressure, so its cache entries are stale and produce
  noisier logit corrections.
\end{itemize}

\subsubsection{FreeTTA EM Weight Distribution}

The EM weight $\alpha_t = 1 - H_\text{clip}(\bx_t)/\log C$ controls how strongly each sample
updates the class mean.
Its distribution (see \texttt{outputs/dynamics\_analysis/uncertainty\_em\_weight.png})
reflects the per-dataset difficulty:

\begin{itemize}
  \item \textbf{Caltech} ($\bar{\alpha}=0.826$): strongly right-skewed toward $\alpha \approx 1.0$;
  CLIP is so confident on Caltech that almost every sample produces a near-maximum-weight M-step.
  Risk of overconfident updates propagating noise.

  \item \textbf{EuroSAT} ($\bar{\alpha}=0.198$): left-skewed toward $\alpha \approx 0$;
  CLIP is uncertain, updates are small.
  FreeTTA compensates with $810$ samples per class (each contributing a small update), yielding
  cumulative drift of $1.14$---equal to Caltech despite $\bar{\alpha}$ being $4\times$ smaller.

  \item \textbf{DTD} ($\bar{\alpha}=0.253$): bimodal distribution; a subset of texture images have
  distinctive visual signatures CLIP recognises confidently, while the rest spread probability
  uniformly. This bimodality explains TDA's advantage: high-$\alpha$ samples dominate the cache,
  making it more informative than FreeTTA's mean (which averages both modes).
\end{itemize}

% ── 3. Distribution Modeling Analysis ────────────────────────────────────────
\subsection{Distribution Modeling Analysis}

\subsubsection{Feature Geometry: Oracle Upper Bounds}

To understand whether the feature space favours local lookups (TDA) or global centroids (FreeTTA),
we compute two oracle classifiers on the \emph{ground-truth} labels:
\begin{enumerate}
  \item \textbf{Oracle Centroid}: assign class $\argmax_c \bmu_c^\ast \cdot \bx$ where
  $\bmu_c^\ast$ is the true class mean of CLIP features.
  \item \textbf{Oracle 1-NN}: assign the label of the nearest neighbour in the test set
  (leave-one-out).
\end{enumerate}

The Geometry Alignment Score $\text{GAS} = \text{OracleCentroid} - \text{Oracle1NN}$ measures
which paradigm the data structure supports.

\begin{table}[H]
\centering
\caption{Oracle upper bounds for each paradigm and GAS. GAS $> 0$ predicts FreeTTA advantage;
GAS $< 0$ predicts TDA advantage.}
\label{tab:oracle_geometry}
\begin{tabular}{@{}lrrrrl@{}}
\toprule
Dataset & Oracle Centroid & Oracle 1-NN & GAS & Cache Pressure & Winner \\
\midrule
Caltech  & 97.44\% & 91.85\% & $+0.056$ & 4.93$\times$   & FreeTTA \\
DTD      & 73.19\% & 63.62\% & $+0.096$ & 8.0$\times$    & TDA (exception) \\
EuroSAT  & 78.26\% & 89.80\% & $-0.115$ & 162.0$\times$  & FreeTTA \\
Pets     & 91.50\% & 83.46\% & $+0.080$ & 19.8$\times$   & FreeTTA \\
ImageNet & 82.72\% & 41.89\% & $+0.408$ & 2.0$\times$    & Tie \\
\bottomrule
\end{tabular}
\end{table}

\textbf{GAS predicts the winner on 4 of 5 datasets.}
EuroSAT's negative GAS ($-0.115$) predicts TDA should win---yet FreeTTA wins by $+6.0\%$.
The reason: EuroSAT's GAS measures the \emph{ground-truth} centroid advantage; FreeTTA's adapted
centroid $\bmu_c$ is far from the ground truth at stream start.
The 162$\times$ cache pressure destroys TDA's memory before its cache can accumulate meaningful
signal, while FreeTTA's global mean---even when imperfect---competes.

DTD is the other exception: GAS $= +0.096$ nominally favours centroid-based methods, yet TDA
wins.
The cause is intra-class texture diversity: DTD's within-class variance is high
($\sigma$-trace $= 0.333$), meaning the true centroid is a blurred average of diverse textures.
TDA's cached \emph{exemplars} preserve this diversity; FreeTTA's mean collapses it.

\subsubsection{Logit-Space PCA}

PCA of the adapted logit vectors (see \texttt{outputs/dynamics\_analysis/distribution\_pca\_*.png})
reveals the structural difference between the two methods' corrections:

\begin{itemize}
  \item \textbf{TDA logits} cluster tightly in PCA space: correct predictions occupy a compact
  region (low PC1 variance) and wrong predictions form a separate sparse cloud.
  This reflects TDA's hard, sparse correction: a single cache hit dominates all logits.

  \item \textbf{FreeTTA logits} show a smooth gradient from correct (high-confidence) to wrong
  (low-confidence) with significant overlap---the soft mean update yields a richer distribution.

  \item \textbf{CLIP logits} are near-uniform and lie close to the origin in logit space on
  all datasets, confirming the calibration failure observed in the entropy analysis.
\end{itemize}

\subsubsection{Distribution Evolution: Covariance Trace and Prior Entropy}

Beyond point estimates, we track two distributional quantities over the stream:

\textbf{$\sigma$-trace} (empirical spread of adapted CLIP features):
In all datasets, $\sigma$-trace decreases over time as FreeTTA's adapted means pull predictions
toward class centroids.
The rate of decrease is fastest on EuroSAT (810 samples/class) and slowest on ImageNet (10 samples/class).

\textbf{Prior entropy $H(\pi)$}: the entropy of FreeTTA's class-prior estimates.
On EuroSAT ($H(\pi)=0.72$) the model concentrates probability on fewer classes as it
discovers the 10-class structure.
On ImageNet ($H(\pi)=3.97 \approx \log 1000 / 1.74$) the prior remains near-uniform:
with only 10 samples per class, FreeTTA cannot identify the 1,000-class structure.

% ── 4. Computational Efficiency ──────────────────────────────────────────────
\subsection{Computational Efficiency}

\subsubsection{Per-Sample Latency}

All experiments were run on CPU.
Table~\ref{tab:efficiency} reports per-sample latency computed from total wall-clock time divided
by stream length.

\begin{table}[H]
\centering
\caption{Computational efficiency on CPU (all times in milliseconds per sample).
Speedup = TDA time / FreeTTA time.}
\label{tab:efficiency}
\begin{tabular}{@{}lrrrrl@{}}
\toprule
Dataset & $N$ & TDA (ms/spl) & FreeTTA (ms/spl) & Speedup & Bottleneck \\
\midrule
Caltech  & 2,465  & 0.394 & 0.294 & 1.34$\times$ & TDA cache lookup  \\
DTD      & 1,880  & 0.375 & 0.187 & 2.01$\times$ & TDA pos+neg cache \\
EuroSAT  & 8,100  & 0.353 & 0.169 & 2.09$\times$ & TDA cache eviction \\
Pets     & 3,669  & 0.381 & 0.192 & 1.99$\times$ & TDA cache growth \\
ImageNet & 10,000 & 1.961 & 0.703 & 2.79$\times$ & TDA $10^3$-class cache \\
\bottomrule
\end{tabular}
\end{table}

FreeTTA is $1.3$--$2.8\times$ faster than TDA across all datasets.
The speedup grows with dataset size and class count, explaining why ImageNet shows the largest
advantage ($2.79\times$).

\subsubsection{Why TDA Is Slower: Algorithmic Analysis}

TDA's per-sample cost is $O(D \cdot S_t)$ where $S_t = |\mathcal{P}| + |\mathcal{N}|$ is the
total number of cache slots occupied at step $t$.
$S_t$ grows monotonically until cache saturation, so TDA becomes \emph{slower} as the stream
progresses:
\begin{equation}
  t_\text{TDA}(t) \propto D \cdot \min(t,\, C \cdot (K_+ + K_-))
\end{equation}
For ImageNet with $C = 1000$, saturation only occurs after $\sim\!20{,}000$ steps---meaning
the entire 10,000-sample stream runs during the unsaturated, growing-cost phase.
At stream end, TDA is doing $D \times (\sim\!2600)$ dot products per sample.

FreeTTA's per-sample cost is $O(D \cdot C)$ \emph{constant}:
the M-step is one vector update plus $C$ comparisons, independent of stream length.

\subsubsection{Memory Usage}

Memory requirements (estimated at end of stream):
\begin{itemize}
  \item \textbf{TDA}: occupied cache slots $\times$ $D$ $\times$ 4 bytes (float32).
  Caltech: $\sim$366\,KB; DTD: $\sim$431\,KB; EuroSAT: $\sim$84\,KB; Pets: $\sim$258\,KB;
  ImageNet: $\sim$5.3\,MB ($2{,}597$ slots $\times 512 \times 4$).

  \item \textbf{FreeTTA}: $C \times D \times 4$ bytes (one mean per class).
  Caltech: 205\,KB; DTD: 96\,KB; EuroSAT: 20\,KB; Pets: 76\,KB; ImageNet: 2.0\,MB.
\end{itemize}

FreeTTA uses $1.5$--$4\times$ less memory than TDA, with the advantage growing with class count
(ImageNet: 2.0\,MB vs 5.3\,MB).
At scale (e.g., 100K classes), TDA's cache would be prohibitive while FreeTTA remains tractable.

\subsubsection{Scalability Conclusion}

Both methods scale linearly with $D$ (fixed at 512 for CLIP ViT-B/16).
TDA scales as $O(D \cdot C \cdot K)$ in memory and $O(D \cdot S_t)$ in time (growing with stream).
FreeTTA scales as $O(D \cdot C)$ in both memory and time (constant per sample).
For open-world classification with thousands of classes, FreeTTA's complexity advantage is decisive.

% ══════════════════════════════════════════════════════════════════════════════
\section{Where Each Algorithm Always Wins and Loses}
% ══════════════════════════════════════════════════════════════════════════════

\subsection{FreeTTA: Conditions for Winning}

\begin{enumerate}
  \item \textbf{High domain shift from internet distribution} (EuroSAT)\\
  \textit{Why}: Text anchor $\bw_c$ is far from the true visual centroid $\bmu_c$.
  Any correction via M-step moves toward the correct distribution.
  \textit{Signal}: CLIP accuracy $< 60\%$ on classes with smooth, learnable visual structure.

  \item \textbf{High samples-per-class ratio} (EuroSAT: 810/class)\\
  \textit{Why}: Mean estimate variance $\sim 1/n_c$; more samples $\to$ better centroid estimate.
  \textit{Signal}: $N/C > 50$.

  \item \textbf{High cache pressure} (EuroSAT: 162$\times$)\\
  \textit{Why}: TDA's memory is overwhelmed; FreeTTA's global mean is unaffected by capacity.

  \item \textbf{Negative or low GAS}\\
  \textit{Why}: When the CLIP anchor is severely wrong, even suboptimal centroid correction beats local lookups.
\end{enumerate}

\subsection{FreeTTA: Conditions for Failing}

\begin{enumerate}
  \item \textbf{Discrete non-overlapping categories where exact cached features outperform global means} (DTD)\\
  \textit{Why}: For texture classification, the within-class visual diversity is large but inter-class boundaries are sharp.
  FreeTTA's global mean blends diverse texture samples into a centroid that may sit between clusters, while TDA's exact cached exemplars remain discriminative.

  \item \textbf{Very few samples per class} (ImageNet: 10/class)\\
  \textit{Why}: With $\eta=0.1$ and 10 updates, $\bmu_c$ shifts by $<10\%$ of the distance to the sample centroid.
  Insufficient to correct anything.

  \item \textbf{Overconfident CLIP on wrong classes} (adversarial scenario)\\
  \textit{Why}: High $\alpha_t$ + high $w_{\text{upd}}$ on wrong prediction $\to$ M-step moves mean in wrong direction, self-reinforcing error.
\end{enumerate}

\subsection{TDA: Conditions for Winning}

\begin{enumerate}
  \item \textbf{Discrete, non-overlapping category structure} (DTD: 47 texture types)\\
  \textit{Why}: When categories are visually distinct and non-overlapping (``bumpy'' vs.\ ``dotted''), cached exemplars transfer reliably to future test samples from the same class.
  Tighter entropy gates ($\theta_{\text{low}}=0.05$, $\theta_{\text{high}}=0.40$) select only high-confidence texture exemplars, producing a high-precision cache.

  \item \textbf{Low cache pressure} (ImageNet: 2.0$\times$; necessary but not sufficient---see Caltech at 4.93$\times$ where TDA still makes more harmful than beneficial flips)\\
  \textit{Why}: Cache remains representative; the most recent confident samples are genuinely informative.

  \item \textbf{Stream locality} (natural order, class-clustered streams)\\
  \textit{Why}: Consecutive same-class samples fill the cache quickly with high-quality examples.
\end{enumerate}

\subsection{TDA: Conditions for Failing}

\begin{enumerate}
  \item \textbf{High cache pressure} (EuroSAT: 162$\times$)\\
  \textit{Why}: Cache slots overwritten so frequently the ``cache'' is effectively the last 3 random samples per class.

  \item \textbf{Cold-start penalty} (EuroSAT break-even at step 941)\\
  \textit{Why}: Empty cache $\to$ zero or noisy adjustment $\to$ early stream is pure CLIP or worse.

  \item \textbf{Confident wrong predictions entering the cache}\\
  \textit{Why}: TDA inserts samples with entropy $< \theta_{\text{low}}$ under predicted label.
  A confidently misclassified sample becomes a poison entry, reinforcing future errors for similar inputs.
\end{enumerate}

% ══════════════════════════════════════════════════════════════════════════════
\section{Future Work and Suggested Improvements}
% ══════════════════════════════════════════════════════════════════════════════

\subsection{Improving TDA}

\subsubsection{Adaptive Cache Capacity}
\textbf{Problem}: Fixed capacity $K=3$ fails at 162$\times$ pressure.\\
\textbf{Solution}: Reservoir sampling---each new sample replaces a random existing entry with probability $1/k$ (number seen so far).
Alternatively, set capacity dynamically:
\begin{equation}
  K(c) = \max\!\left(3,\; \left\lfloor\sqrt{N_{\text{est}} / C}\right\rfloor\right)
\end{equation}
\textbf{Expected gain}: EuroSAT cache pressure drops from 162$\times$ to $\approx$28$\times$.

\subsubsection{Verified Cache Updates}
\textbf{Problem}: Misclassified confident samples poison the cache.\\
\textbf{Solution}: Only cache $\bx_t$ as class $c$ if the majority of existing $\mathcal{P}_c$ entries are more similar to $\bx_t$ than to any other class's cache (self-consistency check).

\subsubsection{Temporal Weighting (Recency Bias)}
\textbf{Problem}: FIFO treats all cached samples equally; old samples may be stale.\\
\textbf{Solution}: Decay cache contributions exponentially:
\begin{equation}
  z^{\text{pos}}_c = \sum_{\mathbf{v}_i \in \mathcal{P}_c} \gamma^{t - t_i} (\bx_t^\top \mathbf{v}_i), \quad \gamma < 1
\end{equation}

\subsection{Improving FreeTTA}

\subsubsection{Per-Class Adaptive Learning Rate}
\textbf{Problem}: Fixed $\eta=0.1$ is too large for rare classes and too small for common classes.\\
\textbf{Solution}: Robbins--Monro schedule:
\begin{equation}
  \eta_c(n_c) = \frac{\eta_0}{1 + \beta \cdot n_c}
\end{equation}
where $n_c$ is the number of M-step updates applied to class $c$.
Guarantees mean-squared convergence to the true centroid under mild conditions.

\subsubsection{Confidence Gating for M-Step}
\textbf{Problem}: Low-confidence M-steps add noise to class means.\\
\textbf{Solution}: Skip the M-step if $p_{\text{adapt}}(\hat{y} \mid \bx_t) < \delta$ (threshold $\delta = 0.5$).
Fewer but higher-quality updates.

\subsubsection{Multi-Template Mean Initialisation}
\textbf{Problem}: Single template ``a photo of $\{\text{class}\}$'' may be a poor starting point.\\
\textbf{Solution}:
\begin{equation}
  \bmu_c^{(0)} = \mathrm{normalize}\!\left(\sum_{k} f_T(t_c^{(k)})\right)
\end{equation}
where $t_c^{(k)}$ are diverse prompts (``a photo of $c$'', ``an image of $c$'', ``$c$ in a photo'', etc.).

\subsubsection{Momentum for Mean Updates}
\textbf{Problem}: Single-sample M-step is noisy; one outlier can shift $\bmu_c$ significantly.\\
\textbf{Solution}: Exponential moving average momentum:
\begin{align}
  \mathbf{m}_c &\leftarrow \beta \mathbf{m}_c + (1-\beta) w_{\text{upd}}(\bx_t - \bmu_c) \\
  \bmu_c &\leftarrow \bmu_c + \eta \mathbf{m}_c
\end{align}
Momentum $\beta=0.9$ smooths gradient estimates.

\subsubsection{Full Covariance Mixture Model (Research Direction)}
\textbf{Problem}: Isotropic Gaussian assumption fails for anisotropic texture classes.\\
\textbf{Solution}: Track diagonal covariance:
\begin{equation}
  \boldsymbol{\Sigma}_c \leftarrow (1-\eta)\boldsymbol{\Sigma}_c + \eta w (\bx - \bmu_c)(\bx - \bmu_c)^\top
\end{equation}
\begin{equation}
  z^{\text{adapt}}_c = -\tfrac{1}{2}(\bx - \bmu_c)^\top \boldsymbol{\Sigma}_c^{-1}(\bx - \bmu_c)
\end{equation}
Computational cost $O(D^2)$ per step per class---expensive at $D=512$, but worth exploring for DTD.

\subsection{Hybrid Approaches}

\subsubsection{Cache-Pressure-Adaptive Ensemble}
\begin{equation}
  \bz^{\text{final}} = \begin{cases}
    \bz^{\text{clip}} + w_g \bz^{\text{FreeTTA}} & \text{if cache\_fullness} < 0.5 \\
    \bz^{\text{clip}} + w_l \bz^{\text{TDA}} + (1-w_l)\bz^{\text{FreeTTA}} & \text{otherwise}
  \end{cases}
\end{equation}

\subsubsection{FreeTTA-Initialised TDA}
Use FreeTTA's adapted mean as virtual initial cache entries to solve TDA's cold-start:
\begin{equation}
  \mathcal{P}_c^{(0)} = \{\bmu_c^{\text{FreeTTA}}\} \quad\text{(synthetic anchor)}
\end{equation}
Real samples gradually replace the virtual anchor.
Expected gain on EuroSAT: cold-start accuracy jumps from 33.8\% to $\approx$56\%.

\subsubsection{Meta-Learned Hyperparameter Prediction}
Train a lightweight meta-model that predicts optimal $(\eta, K, \theta_{\text{low}})$ from label-free dataset statistics:
stream length $N$, estimated class count $C$ (from top-1 entropy distribution), and mean CLIP confidence.
Removes the need for per-dataset manual tuning.

% ══════════════════════════════════════════════════════════════════════════════
\section*{Summary}
\addcontentsline{toc}{section}{Summary}
% ══════════════════════════════════════════════════════════════════════════════

CLIP's zero-shot power comes with a critical limitation: its text anchors are fixed at inference time, making it fragile under domain shift.
Test-Time Adaptation addresses this by using the test stream itself as adaptation signal.
Our comparison of two fundamentally different TTA philosophies---cache-based (TDA) and statistics-based (FreeTTA)---across five diverse datasets yields five main conclusions:

\begin{enumerate}
  \item \textbf{FreeTTA is the better default}: wins or ties on 4 of 5 datasets, with gains correlated directly to domain shift severity.
  \item \textbf{Cache pressure is TDA's fatal flaw}: performance collapses at 162$\times$ pressure (EuroSAT), while FreeTTA's global mean is unaffected by capacity constraints.
  \item \textbf{FreeTTA's adaptation quality is higher}: Beneficial Flip Precision consistently exceeds TDA's, meaning FreeTTA's changes are more reliably improvements.
  \item \textbf{FreeTTA adapts faster on every dataset}: break-even is within the first 0.5\% of the stream on DTD, EuroSAT, Pets, and ImageNet; TDA's cold-start can approach 12\% (EuroSAT step 941), and TDA never reaches net positive on Caltech.
  \item \textbf{TDA's niche is discrete texture classification} (DTD): where intra-class visual diversity is high but inter-class boundaries are sharp, TDA's exact cached exemplars outperform FreeTTA's blended class means.
\end{enumerate}

Both methods are complementary.
A hybrid combining FreeTTA's warm initialisation with TDA's local lookup would address both methods' main failure modes.

\bigskip
\noindent
\textit{Plots referenced in this report are available in \texttt{outputs/comparative\_analysis/}:
\texttt{accuracy\_summary.png}, \texttt{flip\_analysis.png},
\texttt{entropy\_confidence\_summary.png}, \texttt{disagreement\_analysis.png},
\texttt{failure\_bucket\_summary.png}, \texttt{geometry\_analysis.png},
\texttt{latency\_analysis.png}.}

\begin{thebibliography}{9}
\bibitem{radford2021clip}
A.\ Radford et al.,
``Learning Transferable Visual Models From Natural Language Supervision,''
\textit{ICML}, 2021.

\bibitem{tda2023}
J.\ Karmanov et al.,
``Efficient Test-Time Adaptation of Vision-Language Models,''
\textit{CVPR}, 2024.

\bibitem{freetta2022}
Y.\ Su et al.,
``Revisiting Classifier: Diffusing Semantics into Classifier for Generalized Few-Shot Learning,''
\textit{CVPR}, 2022.

\bibitem{eurosat}
P.\ Helber et al.,
``EuroSAT: A Novel Dataset and Deep Learning Benchmark for Land Use and Land Cover Classification,''
\textit{IEEE J.\ Sel.\ Topics Appl.\ Earth Obs.\ Remote Sens.}, 2019.

\bibitem{dtd}
M.\ Cimpoi et al.,
``Describing Textures in the Wild,''
\textit{CVPR}, 2014.
\end{thebibliography}

\end{document}
